% !TeX root = ../../Book5.tex
% 纲要标题：### 13.3 半单分解
% 纲要位置：工作记录/计划纲要.md，第 4155 行。

\section{半单分解}
\label{b5:sec:1303}

% 纲要范围：
%
% * 半单李代数分解为单理想
% * 理想与正交补
% * 分解的唯一性
%

一般表示中的不变子空间未必有不变补。伴随表示在半单情形却有一个特别直接的补：用 Killing 型取正交补。这是半单结构分解的来源，而非任意表示完全可约性的预先使用。

\subsection{半单李代数分解为单理想}
\label{b5:subsec:1303:01}

\begin{theorem}[半单李代数的理想分解]\label{b5:thm:semisimple-ideal-decomposition}
设 \(\mathfrak g\) 是有限维复半单李代数。则存在简单理想
\(\mathfrak g_1,\ldots,\mathfrak g_r\)，使
\[
\mathfrak g=\mathfrak g_1\oplus\cdots\oplus\mathfrak g_r,
\qquad[\mathfrak g_i,\mathfrak g_j]=0\quad(i\ne j).
\]
此外，\(\mathfrak g\) 的每个理想都是其中若干 \(\mathfrak g_i\) 的直和。这些简单理想作为 \(\mathfrak g\) 的子空间由 \(\mathfrak g\) 唯一确定，只允许改变排列顺序。
\end{theorem}
\begin{proof}
先取任意理想 \(I\)。由\cref{b5:prop:invariant-trace-form}，\(I^\perp\) 是理想，所以 \(J=I\cap I^\perp\) 也是理想。对 \(x,y\in J\)，有 \(\kappa_{\mathfrak g}(x,y)=0\)，再由\cref{b5:lem:killing-restriction-ideal}得到 \(\kappa_J=0\)。Cartan 可解性判据说明 \(J\) 可解，半单性迫使 \(J=0\)。由于整个 Killing 型非退化，
\(\dim I^\perp=\dim\mathfrak g-\dim I\)，故
\[
\mathfrak g=I\oplus I^\perp.
\]
两个理想的括号同时属于二者，因而为零。

如果 \(\mathfrak g\ne0\)，选一个极小非零理想 \(I\)。它的任意自身理想都被 \(I\) 保持，又与 \(I^\perp\) 交换，所以也是 \(\mathfrak g\) 的理想。极小性说明 \(I\) 没有非零真理想；它不能交换，否则会是可解理想。因此 \(I\) 简单。对补 \(I^\perp\) 继续此过程，维数严格减小，有限步便得到所述分解。

设 \(J\) 是任意理想，\(x=x_1+\cdots+x_r\in J\)。若 \(x_i\ne0\)，简单代数 \(\mathfrak g_i\) 的中心为零，所以
\([\mathfrak g_i,x_i]\ne0\)。又因为不同分量交换，
\[
[\mathfrak g_i,x]=[\mathfrak g_i,x_i]\subseteq J\cap\mathfrak g_i.
\]
交集是 \(\mathfrak g_i\) 的理想，故简单性迫使 \(\mathfrak g_i\subseteq J\)。于是 \(x\) 的每个非零分量都属于已包含于 \(J\) 的简单理想，从而 \(J\) 就是这些理想的直和。特别地，极小非零理想恰好是所列的 \(\mathfrak g_i\)，唯一性也随之成立。
\end{proof}

\subsection{理想与正交补}
\label{b5:subsec:1303:02}

\begin{corollary}\label{b5:cor:semisimple-ideal-complement}
设 \(\mathfrak g\) 是有限维复半单李代数，\(I\) 是理想。则
\[
\mathfrak g=I\oplus I^\perp,\qquad
I^\perp=Z_{\mathfrak g}(I),\qquad
\mathfrak g/I\cong I^\perp.
\]
这里正交补取自 Killing 型，中心化子为
\(Z_{\mathfrak g}(I)=\Bigl\{\,x\mid[x,I]=0\,\Bigr\}\)。
而且 \(I\) 与 \(\mathfrak g/I\) 都半单。
\end{corollary}
\begin{proof}
由\cref{b5:thm:semisimple-ideal-decomposition}，\(I\) 是部分简单理想的和，其余分量组成 \(I^\perp\)。后者显然中心化 \(I\)；若一个元素在 \(I\) 分量上的投影也中心化 \(I\)，该投影属于 \(Z(I)=0\)，所以中心化子没有更多元素。商映射在 \(I^\perp\) 上单射且满射，给出同构。两者都由简单理想直和组成，故半单。
\end{proof}

\begin{corollary}\label{b5:cor:semisimple-perfect}
有限维复半单李代数 \(\mathfrak g\) 满足
\[
[\mathfrak g,\mathfrak g]=\mathfrak g,\qquad Z(\mathfrak g)=0.
\]
因此其一维表示全为平凡表示。
\end{corollary}
\begin{proof}
每个非交换简单理想的导出代数是非零理想，因而等于自身；中心是可解理想，因而为零。对\cref{b5:thm:semisimple-ideal-decomposition}中的直和逐项应用即可。一维表示的像交换，必零化导出代数，故零化整个 \(\mathfrak g\)。
\end{proof}

这个正交补只针对理想。一般子代数 \(\mathfrak h\) 的正交补不一定是理想，\(\kappa|_{\mathfrak h}\) 也不一定非退化；例如 \(\mathfrak{sl}_2\) 中的一维子代数 \(\mathbb Ce\) 就是各向同性的。

\subsection{分解的唯一性}
\label{b5:subsec:1303:03}

\begin{example}\label{b5:exm:diagonal-simple-not-ideal}
设 \(\mathfrak s\) 是非零复简单李代数。对角子代数
\[
\Delta\mathfrak s=\Bigl\{\,(x,x)\mid x\in\mathfrak s\,\Bigr\}
\subseteq\mathfrak s\oplus\mathfrak s
\]
同构于 \(\mathfrak s\)，但不是理想。取 \([a,x]\ne0\)，则
\[
\Bigl[(a,0),(x,x)\Bigr]=\Bigl([a,x],0\Bigr)\notin\Delta\mathfrak s.
\]
所以简单子代数可以很多，而理想分解中的简单理想仍只有两个。
\end{example}

分解唯一性说的是具体理想的集合唯一，比仅仅列出抽象同构类型更强。自同构可以置换同构的简单理想，但不可能把一个简单理想送到上述对角子代数，因为自同构保持理想性。这也说明分解中出现重数时仍没有连续选择补空间的自由。
